%% ICAIF 2026 submission -- protocol-first register.
%% CFP (verified 2026-07-16, icaif2026.org): sigconf + anonymous, EIGHT (8)
%% pages TOTAL including figures and references; deadline 2026-08-02 AoE;
%% double-blind. NOTE: the limit INCLUDES references, so the reference list
%% is trimmed and the body targets ~7.25pp.
\documentclass[sigconf,anonymous]{acmart}

\settopmatter{printacmref=false}
\setcopyright{none}
\acmConference[ICAIF '26]{7th ACM International Conference on AI in
Finance}{November 14--17, 2026}{Milan, Italy}
\acmYear{2026}

\usepackage{booktabs}
\usepackage{amsmath}
\usepackage{array}

% Visible TODO marker (kept deliberately loud for the pre-submission pass)
\newcommand{\todo}[1]{\textcolor{red}{[TODO: #1]}}

\newcommand{\G}{\hat\Gamma}
\newcommand{\regret}{\rho}

\begin{document}

\title[Should You Pay to Remember?]{Should You Pay to Remember? A Two-Gate,
Pre-Registered Protocol for Deciding When Regime Retention Pays in Trading
Systems}

\author{Anonymous Author(s)}
\affiliation{%
  \institution{Anonymous Institution}
  \city{}
  \country{}}
\email{anonymous@example.com}

\begin{abstract}
Regime-aware trading systems carry an evaluation problem that is worse than
overfitting: their evaluations cannot fail. A regime-conditioned learner
happily trains, ``specializes,'' and beats a weak baseline on structureless
labels, so a backtest win licenses nothing. We propose and field-test a
pre-deployment protocol built around one number: the \emph{retention
deficit} $\G$, the post-reactivation decision-regret difference between the
deployed always-retrain learner and an identical twin whose per-regime
parameters are protected during dormancy---a two-arm counterfactual any
desk can run before buying retention machinery. The protocol pairs $\G$
with a structure screen (regime-conditioning versus block-shuffled labels
on the decision metric), which its own record demotes: across our field
tests the screen's verdict dissociated from realized value three times,
including its strongest pass coinciding with an inverted outcome. The field
record spans daily crypto and commodities (screen fails, evaluations go
flat, as required), 4-hour crypto (screen passes, $\G=0$, flat again),
and daily U.S. industry portfolios 1926--2025, where $\G$'s sign tracked
the outcome direction in both directions: a positive deficit
($+0.00089\pm0.00011$) coincides with the full pre-registered arm ordering,
a negative one with an inversion, and a frozen-specification replication on
a withheld 64-year era scores 5 of 6 cells. Crisis-window transaction costs
amplify rather than kill the deficit ($+0.00126\to+0.00230$ as costs sweep
25--100\,bps). An externally lodged, announcement-lagged NBER-recession
test---registered with the OSF before recession dates ever touched the
panel---resolved as a hit, with its two-recession concentration disclosed.
A break-even rule converts $\G$ into a dollar-denominated carrying-cost
threshold. All claims are regret and retention claims, not alpha claims; a
live temporal test (diagnostics frozen 2026-06-30, scored 2027-04-07) and a
constituent-level replication are lodged in advance.
\end{abstract}

\begin{CCSXML}
<ccs2012>
<concept>
<concept_id>10010405.10010481.10010487</concept_id>
<concept_desc>Applied computing~Economics</concept_desc>
<concept_significance>500</concept_significance>
</concept>
<concept>
<concept_id>10010147.10010257</concept_id>
<concept_desc>Computing methodologies~Machine learning</concept_desc>
<concept_significance>300</concept_significance>
</concept>
</ccs2012>
\end{CCSXML}
\ccsdesc[500]{Applied computing~Economics}
\ccsdesc[300]{Computing methodologies~Machine learning}

\keywords{regime switching, continual learning, decision-focused learning,
evaluation protocols, pre-registration, transaction costs}

\maketitle

%=============================================================================
\section{Introduction: the evaluation that cannot fail}
\label{sec:intro}
%=============================================================================

Every multi-strategy desk faces the same quiet question: the crisis
playbook has been idle for six years---who is paying to keep it
calibrated, and is that money buying anything? The machine-learning
literature now offers an inventory of retention machinery to answer
``yes'': per-regime expert libraries, replay buffers, parameter isolation,
continual-learning regularizers ported to finance
\cite{philps2018continual,pan2026recap,kirkpatrick2017overcoming,
rolnick2019experience,parisi2019continual}. What it does not offer is a
way to find out that the answer is ``no.''

The reason is structural. A regime-aware evaluation has more researcher
degrees of freedom than an ordinary backtest---the labeler, its
thresholds, the dormancy definition, the evaluation windows---and its
machinery produces organized-looking behavior on any labels at all: a
regime-conditioned learner trains, converges, and ``specializes'' whether
or not the labels carry information. Beating a generic baseline is
therefore not evidence. \emph{A method whose evaluation cannot go flat on
structureless data should not be trusted when it goes sharp.} This paper
is organized around engineering evaluations that can go flat, running
them on real substrates, and reporting the flats, the inversions, and the
hits in one ledger.

We work in the register of decision regret, not P\&L: the per-day gap
between the portfolio chosen from a model's predictions and the portfolio
an oracle would have chosen, the predict-then-optimize
lesson that prediction error and decision quality dissociate
\cite{elmachtoub2022smart,mandi2024decision}. The question the protocol
prices is narrow and universal: when a dormant market regime returns,
does a capacity-bounded learner that spent the dormancy retraining on
whatever was active pay a measurable penalty in the first days of the
reactivation---the days when decisions matter most---relative to a twin
that kept the dormant regime's parameters untouched? If yes, retention is
worth money and the break-even arithmetic says how much; if no, every
dollar of retention machinery is premium paid on a claim that will never
be filed. Recent theory makes the same point from first principles: a
bounded agent should not retain what it can cheaply relearn
\cite{kumar2025continual}; our contribution is the measurement that
decides the question per substrate, before deployment.

\textbf{Contributions.} (1)~\emph{A two-gate pre-deployment protocol} for
regime-aware trading systems: a cheap structure screen run first, a
binding deficit diagnostic run second, and a break-even carrying-cost
rule that converts the diagnostic into dollars
(Section~\ref{sec:protocol}). (2)~\emph{The diagnostic itself}: $\G$, a
two-arm counterfactual (deployed learner versus its dormancy-protected
twin) measured in post-reactivation decision regret---one number, causal
by construction, cheap enough to pre-register and run on any book.
(3)~\emph{A resolved field record} (Section~\ref{sec:record}): seven
substrates spanning crypto, commodities, and a century of daily U.S.
equity industry data, every null and inversion included, with the
screen's honest demotion, a costs battery in which the deficit
strengthens net of crisis-window costs (Section~\ref{sec:costs}), an
externally custodied forward test that has since resolved, and a live
temporal test whose diagnostics are already frozen
(Section~\ref{sec:signrule}). The retention mechanism, its theory, and
the controlled dissociation studies are developed in a companion
manuscript \cite{companion2026}; this paper is the
protocol, the number, and the record a practitioner needs to decide
whether to pay to remember.

\subsection{Position among existing practice}
\label{sec:related}

Three literatures meet at this question and each leaves it unpriced.
\emph{Regime-switching finance}---Markov-switching models
\cite{hamilton1989new}, regime effects in returns \cite{ang2012regime},
regime-conditioned allocation \cite{guidolin2007asset,
nystrup2018dynamic,nystrup2020greedy}, regime-weighted forecast
combination \cite{elliott2005optimal}---conditions a model or its
combination weights on an inferred state, implicitly assuming every
component remains fully estimated while its state is dormant. That
assumption is exactly what $\G$ measures the failure of under bounded
capacity. \emph{Continual learning in finance} makes the opposite
unconditional bet: retention machinery is deployed because forgetting
is presumed harmful \cite{philps2018continual,pan2026recap,
kirkpatrick2017overcoming,rolnick2019experience,parisi2019continual},
with evaluation by averaged forgetting metrics
\cite{lopez2017gradient} rather than by a pre-deployment test of
whether the substrate offers anything to retain. Our record contains
real substrates where that bet loses ($\G < 0$ with a realized
inversion), which is why the protocol runs \emph{before} the purchase.
\emph{The backtest-overfitting protocol literature}---deflated Sharpe
ratios \cite{bailey2014deflated}, SPA and Reality-Check family tests
\cite{hansen2005spa,white2000reality}, the broader discipline of
\cite{harvey2016backtesting,lopez2018advances}---is our closest
methodological ancestor: it converts an unfalsifiable claim (``my
strategy works'') into a testable one by charging for selection. The
two-gate protocol does the same for a different claim (``my system
needs to remember''), and inherits that literature's stance: the
deliverable is the test, and the test must be able to say no.

%=============================================================================
\section{The protocol}
\label{sec:protocol}
%=============================================================================

\subsection{Decision regret and probe windows}

A stream of trading periods presents per-asset features $X_t$ and
next-period returns $y_t$, with a causal regime label $r_t$ computed from
information through $t-1$ only. The decision layer is a
cardinality-constrained long-only portfolio (top-$k$ by predicted return,
position cap $w_{\max}$), and the per-period \emph{decision regret} of a
prediction $\hat y$ is
$\regret(\hat y; y) = y^\top z^*(y) - y^\top z(\hat y) \ge 0$: the return
given up relative to the oracle portfolio. Regret, not prediction error,
is the protocol's metric throughout---errors that do not change the
selected book are free, and errors at the selection boundary are not
\cite{elmachtoub2022smart}.

Regimes recur with dormancy: the crisis label may sleep for years. The
protocol's unit of account is the \emph{probe window}: the first
$N_p = 15$ trading days of every regime block whose regime was dormant at
least 90 days. Post-reactivation regret---mean daily regret inside probe
windows---is the primary metric because the reactivation transient is
where retention could pay and where a fresh learner is most exposed;
regime-conditioned finance has long known returns behave differently
across states \cite{hamilton1989new,ang2012regime,guidolin2007asset,
nystrup2018dynamic}, but conditions a model on the state without asking
who maintained the model through the state's absence
\cite{elliott2005optimal}.

\textbf{Conventions, fixed once.} Every comparison in this paper is
walk-forward with causal features; every claim is a pre-registered
ordering prediction tested across seeds (100 where stated, else 20)
with Welch tests and Holm--Bonferroni correction inside each
registered pair family; paired designs expose every arm to identical
data streams per seed. Where a design's seeds vary only
initialization---the walk-forward batteries---the resulting intervals
are implementation noise conditional on one market history, and every
table says so. Refuted predictions are reported at the same prominence
as confirmed ones; this paper contains both.

\subsection{The structure screen, honestly demoted}
\label{sec:screen}

\textbf{Screen (structure).} Fit one model per regime label, walk-forward;
compare its decision regret against the identical machinery run on
block-shuffled labels ($\ge 50$ draws), summarized as a $z$-score against
the shuffle distribution (the empirical permutation $p$ cannot fall below
$1/51$; a normal fit to the control distribution gives the indicative
$p$), with a split-half stability check and a comparison against a pooled
unconditioned model. Block-shuffling preserves label persistence while
destroying the label-return association, so the control inherits the
machinery's ability to look organized on nothing. The screen asks only:
\emph{do these labels carry decision-relevant information at all?} It
costs minutes and requires no retention machinery.

We designed the screen as a gate---no pass, no claims---and its own field
record demoted it. Across the substrates of Section~\ref{sec:record} the
screen's verdict dissociated from realized retention value three times:
a passing screen preceded an inverted outcome (daily equities under a
volatility-band labeler), the strongest pass we ever measured
($z = 23.5$, withheld 1926--1989 era) coincided with an inversion, and a
failing screen preceded a positive, cost-robust deficit (the NBER
labeler). The screen therefore \emph{licenses nothing}. It survives in
the protocol for its original falsification role---where it fails,
a null evaluation is the predicted outcome and a sharp one is suspect
---and as a cheap negative control: machinery that ``wins'' where the
screen fails is fitting noise.

\subsection{The deficit \texorpdfstring{$\G$}{Gamma}: one number from two
arms}
\label{sec:gamma}

\textbf{Diagnostic (binding).} Run two arms of identical capacity,
features, and learning rates over the same walk-forward stream:
\begin{itemize}
  \item \textbf{A1 (deployed default)}: a capacity-bounded learner that
  always trains on the currently active regime---the honest model of any
  allocation policy that follows current reward, whether a retraining
  schedule, a learned router, or recent-performance capital weighting;
  \item \textbf{A9 (protected twin)}: the same learner with per-regime
  parameters pinned---each regime's parameters are trained only on that
  regime's data and are untouched during its dormancy.
\end{itemize}
The \emph{retention deficit} is
$\G = \overline{\regret}_{\text{A1}} - \overline{\regret}_{\text{A9}}$
over probe windows: what dormancy-period retraining actually cost, in
regret per probe day. In continual-learning terms $\G$ is a
decision-regret backward-transfer measure against a multi-head oracle
\cite{lopez2017gradient}; in econometric terms it descends from the
post-break estimation-window question---how much pre-break knowledge
should a forecaster retain \cite{pesaran2007selection}. It is causal by
construction (the arms differ in exactly one design bit), cheap (two
arms, no mechanism), and signed: $\G \le 0$ is a finding, not a failure.
Where $\G$ is positive, any isolation-preserving implementation
collects it;
where it is not, no retention mechanism can be bought, because there is
nothing to collect.

Both arms are stand-ins for classes, which is what makes the two-arm
run sufficient. A1 represents every allocation policy whose training
resource follows current reward---a retraining schedule, a learned
router over experts, recent-performance capital weighting across
desks---all of which starve the dormant regime identically, because a
dormant regime emits no reward to follow. A9 represents every
\emph{isolation-preserving} implementation: a siloed desk, a frozen
per-regime adapter store, an expert library with pinned members---and,
at the cheap end, data retention alone: in the pre-registered replay
control on this panel, a monolith that rehearses stored crisis
episodes collects roughly 99\% of the deficit. The
protocol prices the class boundary, not any implementation; choosing
an implementation is a separate decision the companion studies
\cite{companion2026}. Two practical notes. First, $\G$ is measured
net of \emph{rehearsals}: if near-miss episodes keep the dormant
specialist warm between major events, that recalibration is in the
measurement, which is the economically relevant accounting. Second,
$\G$ should be re-measured net of transaction costs before pricing
(Section~\ref{sec:costs}); in our record costs move the number, in
one case across zero.

\subsection{Pricing the deficit: the break-even rule}
\label{sec:breakeven}

$\G$ is in return units per probe day, so it converts directly to
dollars. Charging a carrying cost $c$ (dollars per idle trading day:
infrastructure, data, attention, the calibration operation) against a
sleeve of notional value $V$, retention of a regime whose dormancy is
$D$ trading days pays over one dormancy--reactivation cycle iff

\begin{center}
\fbox{\parbox{0.88\columnwidth}{\centering
\vspace{2pt}
$c \cdot D \;\le\; \G_{\mathrm{net}} \cdot N_p \cdot V$
\vspace{2pt}
}}
\end{center}

\noindent where $\G_{\mathrm{net}}$ is the deficit measured net of
transaction costs (Section~\ref{sec:costs}) and $N_p$ is the probe
length. The rule is a special case of the break-even analysis developed
in the companion \cite{companion2026}; here it is used only as
arithmetic.

\textbf{Worked example (a \$100M crisis sleeve).} Take the measured
net-of-cost deficit at the 25\,bps crisis-cost tier,
$\G_{\mathrm{net}} = +0.00126$ per probe day
(Section~\ref{sec:costs}): on \$100M that is \$126{,}000 per probe day,
or ${\approx}\$1.89$M of avoidable regret per reactivation
($N_p = 15$). The June-2020 recession-cell reactivation in our record
arrived after $D = 2{,}445$ trading days of dormancy; the rule prices
retention as paying iff the sleeve's carrying cost is below
$\$1.89\text{M}/2{,}445 \approx \$770$ per idle trading day---about
\$195{,}000 per year. A fully staffed idle desk at \$1M/year does not
pay on this deficit; a frozen parameter store with a nightly calibration
job at \$50K/year pays roughly fourfold. The rule's two-sided honesty is
the point: it prices some retention out, and that is the protocol
working.

%=============================================================================
\section{The field record}
\label{sec:record}
%=============================================================================

Table~\ref{tab:master} is the paper's spine: every substrate the
protocol has touched, with the screen verdict, the measured deficit, and
what the ten-arm evaluation machine (the companion's arm family,
spanning reward-driven and retaining allocations \cite{companion2026})
actually did. Conventions: all labels causal; all comparisons
walk-forward; Welch tests with Holm correction inside each
pre-registered family; 100 seeds where stated, else 20. One caveat
governs every real-data row and is repeated in
Section~\ref{sec:limits}: each walk-forward window is a single market
history, so intervals are seed noise conditional on that history, not
sampling uncertainty across markets.

\begin{table*}[t]
\centering
\caption{The field record: substrate $\times$ (structure screen, deficit
$\G$, evaluation outcome). $\G$ is in decision-regret units per probe
day (A1$-$A9, probe 15 days, dormancy $\ge 90$ days; $\pm$ is a 95\% CI
over seeds, conditional on each window's single history). ``Ordering''
= the pre-registered arm ordering (retaining arms ahead of reward-driven
arms); ``inversion'' = the deployed always-retrain baseline wins;
``predates gate~2'' = the battery ran before the deficit diagnostic
entered the protocol. Every row's result files are released.}
\label{tab:master}
\footnotesize
\begin{tabular}{@{}>{\raggedright\arraybackslash}p{148pt}>{\raggedright\arraybackslash}p{94pt}>{\raggedright\arraybackslash}p{96pt}>{\raggedright\arraybackslash}p{136pt}@{}}
\toprule
\textbf{Substrate (span; labeler)} & \textbf{Structure screen} &
\textbf{Deficit $\G$ (A1$-$A9)} & \textbf{What the machine did} \\
\midrule
Crypto, daily (5 pairs, 2021--25; vol/trend + filtered-vol) &
fail ($z=-0.66$ / $-0.17$) & --- (predates gate 2) &
flat, as predicted (min Holm $p=0.17$; 100 seeds) \\
Commodities, daily (3 series, 2015--25; both labelers) &
fail ($z=-0.14$ / $+0.16$) & --- (predates gate 2) &
no dissociation run; claims forbidden by screen \\
Crypto, 4-hour (bar-native filtered-vol; confirmed cell) &
pass ($+3.15\%$, $z=2.32$) & $+0.000002$ (zero to 4 dp) &
flat ($p>0.57$ all pairs): structure, nothing to retain \\
\addlinespace
French 49 industries, daily (1990--2025; L1 vol-band) &
pass ($+1.14\%$, $z=3.18$) & $-0.00013 \pm 0.00007$ (100 s) &
\emph{inversion}: A1 beats pool (Holm $1.2{\times}10^{-26}$); reverses
net of costs \\
French 49 industries, daily (1990--2025; L3 drawdown@15\%) &
fail by pooled letter ($z=2.24$) & $+0.00089 \pm 0.00011$ (100 s) &
full ordering (headline pair Holm $5.2{\times}10^{-65}$; binding chain
pair $1.6{\times}10^{-15}$); cost-amplified \\
Withheld era: 36 industries (1926--89; frozen L3 spec, 6 cells) &
strongest pass ($z=23.5$) & sign varies by cell &
sign rule 5/6, both directions; screen's best pass $\to$ inversion \\
NBER recession $\times$ trend, announce-lagged (1990--2025) &
fail ($z=1.74$) & $+0.00081 \pm 0.00033$ &
ordering (Holm $1.6{\times}10^{-4}$); two-recession concentration
disclosed \\
\bottomrule
\end{tabular}
\end{table*}

\textbf{Daily crypto and commodities: the protocol's floor.} Neither
small daily panel shows decision-metric structure under either causal
labeler (all $|z| < 0.7$): regime-conditioned regret is statistically
identical to block-shuffled-label regret. The ten-arm evaluation run on
regime-stitched crypto schedules goes flat exactly as the failed screen
predicts (no pair survives Holm correction; min Holm $p = 0.17$), and
the numerically best arm is the one that \emph{learns nothing}---a
fixed-strategy combination---which sharpens the null: where labels carry
no structure, machinery that adapts to them buys less than machinery
that ignores them. This is the ``can go flat'' property demonstrated,
not asserted.

\textbf{4-hour crypto: structure without a deficit.} A richer intraday
feature inventory locates the program's first screen pass (951 episodes;
conditioning beats shuffled labels by $3.15\%$, $z = 2.32$, both halves
positive)---with its selection caveats stated at the same prominence:
it is one confirmed cell from a ten-cell screen, and the confirmation
reused the screen's data span, so the gap estimate is biased upward by
selection on the maximum. The evaluation is flat anyway---and the
two-arm counterfactual explains why retrospectively: $\G = +0.000002$,
zero to four decimals. Reactivations are elevated for every arm
including the protected twin; regime onsets are intrinsically hard
here, but nothing is forgotten---at this bar density a fresh learner
refits within the probe window. We are explicit about the epistemic
order: the pre-registered prediction for this cell was separation, it
was refuted, and the zero deficit is the diagnosis made after the
refutation; the diagnosis's honest test is every substrate scored
since. This cell is why the deficit, not the screen, is the binding
diagnostic: structure is necessary for anything to be learnable, and
irrelevant to whether retention pays.

\textbf{Daily equities, 1990--2025: the deficit exists, and its sign
prices the pool.} On the Ken French 49 industry portfolios (daily,
value-weighted, 1990--2025; 9{,}067 days), with designs and predictions
pre-registered before each run, two labelers deliver the two halves of
the story. Under the volatility-band labeler (L1) the screen passes and
the deficit is \emph{negative} ($-0.00013 \pm 0.00007$): the machine
run anyway produces a significant inversion---the always-retrain
baseline beats the retaining arms in the probe window (Holm
$1.2{\times}10^{-26}$), retention premium paid with no claim to
collect. Under an index-drawdown labeler (L3, 15\% threshold fixed in
the registration before any run) the screen \emph{fails} by the letter
of its pooled sub-criterion while the deficit is positive---the first
positive $\G$ on any real substrate in this program
($+0.00089 \pm 0.00011$, 100 seeds)---and the full pre-registered arm
ordering emerges (headline contrast, retain-vs-retrain, at Holm
$5.2{\times}10^{-65}$; the ordering's binding pair, pool-with-invariance
vs.\ pool-without, at $1.6{\times}10^{-15}$), with the retaining
arms also ahead on overall regret. One labeler-anatomy fact matters for
reading the number: the \emph{union} of the drawdown labeler's crisis
states recurs every one to three years (the 2010, 2011, 2015--16, and
2018 near-misses are all in the inventory), while the specific
crisis-with-uptrend cell sleeps up to 2{,}953 trading days---so
retained parameters do receive live-fire recalibration between major
crises, and $\G$ is measured net of those rehearsals, as a desk's
economics would be. The deficit's evaluated inventory is
27 reactivations; under leave-one-reactivation-out it survives every
single-event exclusion (range $[0.00071, 0.00101]$, all significant),
survives excluding all of calendar-2020 ($+0.00080 \pm 0.00019$), and
is independently positive within the 2010s ($+0.00077 \pm 0.00022$,
$n{=}14$) and 2020s ($+0.00094 \pm 0.00040$, $n{=}12$).

Three disclosures belong beside that result. First, \emph{threshold
fragility}: the deficit lives at essentially one drawdown threshold per
era---at 10/12/20\% on 1990--2025 it is null or negative
($+0.00013 \pm 0.00027$, $-0.00010 \pm 0.00017$,
$-0.00045 \pm 0.00026$)---and the pre-registered robustness prediction
was refuted. Second, a \emph{favorable wrinkle we flag because it was
not predicted}: at 100 seeds the schedule-resampled (stitched)
counterparts, previously flat, turn marginally positive in $\G$ under
both labelers ($+0.00023 \pm 0.00016$ and $+0.00023 \pm 0.00015$)
while their arm families stay flat---consistent either with
single-history optimism in the walk-forward precision or with
stitching destroying exactly the multi-year dormancy structure that
makes retention pay; the lodged constituent-level test discriminates.
Third, \emph{mechanism honesty}: registered probes found
the deficit does not decay across the probe window and follows no
dormancy or rehearsal covariate; it behaves as a persistent
regime-conditional allocation deficit---the reward-following learner
serves crisis windows persistently worse, the learner-side analogue of
slow-moving capital \cite{duffie2010presidential,mitchell2007slow}---not
demonstrated eviction-forgetting, which remains a controlled-setting
result in the companion \cite{companion2026}. The deployment
consequence (the counterfactual prices retention either way) is
unchanged; the mechanism attribution is open, and we say so.

{\sloppy
\textbf{The withheld era, 1926--1989.} The byte-identical frozen
specification---same features, arms, tests, and the lodged sign
hypothesis; zero new researcher degrees of freedom---was replicated on
36 industries over 16{,}983 days the design had never seen, in the
pre-discovery-sample tradition \cite{linnainmaa2018history,
goyal2008comprehensive}. The screen delivers its strongest pass ever
($z = 23.5$) while the frozen 15\% cell \emph{inverts}
($\G = -0.00024 \pm 0.00008$; the always-retrain baseline wins, Holm
$2.1{\times}10^{-3}$)---the dissociation that finally demoted the
screen. The deficit's sign, by contrast, scored: of six withheld cells
(thresholds 10/12/15/20\%; two sub-eras at 15\%), five are consistent
with the lodged sign rule, in both directions, including two
positive-deficit cells reproducing the full ordering out-of-sample
(10\%: $\G = +0.00020 \pm 0.00009$, Holm $1.9{\times}10^{-6}$;
1958--89: $\G = +0.00021 \pm 0.00015$, Holm $2.1{\times}10^{-4}$) and
two negative-deficit cells with the predicted inversion. The single
miss (12\%: $\G$ negative with a CI excluding zero, arm table flat) is scored against
the rule under its strong form and for it under its weak form;
Section~\ref{sec:signrule} carries the full accounting. The paying
threshold \emph{relocated} across eras ($15\% \to 10\%$) as crisis
density shifted---the fragility replicates in pattern, not in location,
so any deployment must estimate the window's location for its own era.
\par}

\textbf{The NBER cell: temporal custody.} The one cell whose ex-ante
status a third party can certify used a labeler with no free threshold
anywhere: regime = NBER recession state $\times$ 50-day trend, the
causal primary variant admitting a recession only from its
\emph{announcement} date. The registration was deposited with the Open
Science Framework at 22:47 ET on 2026-07-15; NBER dates were first
joined to the panel at 22:55 ET. Recession-cell dormancies are
crisis-scale by construction (up to 2{,}557 trading days). The screen
\emph{fails} both variants ($z = 1.74$ primary)---the third
dissociation of screen from value---while the deficit prices positive:
$\G = +0.00081 \pm 0.00033$, strengthening to $+0.00105 \pm 0.00033$
net of 25\,bps crisis-window costs, with the ordering exactly as the
rule demands (Holm $1.6{\times}10^{-4}$ gross; the full retaining chain
at Holm $\le 3.1{\times}10^{-5}$ net). Disclosures at equal prominence:
the lodged leave-one-out and era-blocked slices show the deficit
survives every one of 20 single-event exclusions and all calendar
exclusions (drop-2020: $+0.00042 \pm 0.00026$; drop-2008/09:
$+0.00036 \pm 0.00030$), but its mass concentrates in the two
deep-dormancy recession entries (January 2009, dormancy 1{,}378 days;
June 2020, dormancy 2{,}445 days; dropping the former alone cuts $\G$
to 41\% of headline), and the 2010s era block---which contains zero
recession-cell reactivations---is negative-significant. The cell is
event-robust in the strict sense, two-recession-concentrated, and not
era-consistent; it is materially weaker than the drawdown cell's
robustness profile, and we say so rather than average it away.

%=============================================================================
\section{Costs and economics}
\label{sec:costs}
%=============================================================================

A retention effect that dies under transaction costs is a curiosity;
this one strengthens. The pre-registered cost battery charges
crisis-window effective costs of 25/50/100\,bps under a both-pay
convention (every arm, including the protected twin, charged its own
turnover; an arm-only variant was lodged as sensitivity and matches to
the digit, because the twin's cost term cancels in the regret
difference).

\textbf{The deficit is cost-amplified.} On the positive-deficit equity
cell, the net-of-cost deficit rises monotonically with the cost tier:
$\G_{\mathrm{net}} = +0.00126 \pm 0.00020$ at 25\,bps,
$+0.00161 \pm 0.00021$ at 50, $+0.00230 \pm 0.00024$ at 100, with the
four-arm ordering intact at every tier and the headline contrast
strengthening from Holm $4.6{\times}10^{-19}$ to
$1.3{\times}10^{-29}$. The mechanism is turnover, and it is the
economically interpretable one: a reactivated always-retrain learner
churns its top-$k$ book hardest exactly when spreads are widest
(measured turnover 1.03/day for the deployed baseline versus 0.62 for
the calmest retaining arm). Paying to remember buys, among other
things, \emph{not trading like a fresh learner during a crisis onset}.
The same signature reproduced ex ante on the NBER cell
($+0.00081 \to +0.00105$ net of 25\,bps; turnover 1.04 vs.\ 0.63)---a
labeler the cost battery never saw.

\textbf{The inversion was a gross-only phenomenon.} Under the
volatility-band labeler, where $\G < 0$ and the always-retrain baseline
won gross, the inversion \emph{reverses} net of costs at every tier
(paired difference $-0.00075 \pm 0.00010$ at 25\,bps,
$p = 1.2{\times}10^{-16}$): net of realistic crisis-window costs the
retaining arm wins even where the parameter-level deficit is negative,
because the turnover discipline alone is worth more than the deficit
costs. We report this as economics, not vindication: the sign rule
prices the \emph{parameter} claim, and turnover discipline is in
principle separately purchasable. One sensitivity is
disclosed and not acted on: charging costs inside the structure screen
would flip the drawdown labeler's pooled sub-criterion from
$-0.10\%$ to $+3.0\%$; the screen as registered is a gross-metric
criterion and remains a fail.

\textbf{The desk countermeasure was pre-registered and run: a banded
baseline does not buy the retaining arms' edge.} A no-trade hysteresis
band on the deployed baseline---update the served book only when more
than $b$ names would change, $b \in \{1, 2\}$---is the cheapest desk
implementation of turnover discipline, and its control battery was
lodged with an adverse branch stated before the run (had banding
closed the net gap, the net-superiority claim would have been
withdrawn in favor of ``any turnover-disciplined implementation
collects it,'' at full prominence). The adverse branch is refuted.
Banding is free---gross regret unharmed, net improved---yet the
retaining arm's net advantage stands on both equity cells: on the
positive-deficit (drawdown) cell it beats both banded baselines at
Holm $\le 7.1{\times}10^{-17}$ with $97\%$/$79\%$ of the net gap
intact at $b{=}1$/$b{=}2$, and on the volatility-band cell a
significant residual survives ($+0.00028 \pm 0.00011$, Holm
$2.1{\times}10^{-5}$), with the disclosure that the $b{=}2$ band
closes $62\%$ of that cell's reversal magnitude---the reversal keeps
its label with the attenuation reported alongside. One lodged
magnitude failed and is disclosed rather than renarrated: the bands
cut the baseline's turnover only $4.1\%$/$15.7\%$ against a
registered ${\ge}30\%$, because the crisis churn comes from daily
re-rankings deeper than two names, which defeat narrow bands. The
practitioner conclusion is the battery's point: the net edge is
\emph{selection quality in crisis windows}, not just churn discipline,
and it cannot be bought with a band.

\textbf{What the cost model omits} is market impact, borrow, and
capacity; every substrate outside the two equity batteries remains
gross; and none of these numbers is a P\&L claim
(Section~\ref{sec:limits}). For any future alpha claim the lodged
protocol of record is deflated Sharpe \cite{bailey2014deflated} plus an
SPA/Reality-Check family test \cite{hansen2005spa,white2000reality},
per the backtest-overfitting literature
\cite{harvey2016backtesting,lopez2018advances}.

%=============================================================================
\section{The sign rule: a lodged, once-resolved hypothesis}
\label{sec:signrule}
%=============================================================================

The record above compresses to one deployment rule: \emph{the sign of
$\G$ prices retention}. Positive deficit $\Rightarrow$ the retaining
arms collect it; zero or negative $\Rightarrow$ retention pays premium
without claims. This section states exactly how much evidence that rule
has, in the register the evidence deserves.

\textbf{The confession first.} When first formulated, the rule was
postdictive---written after seeing both 1990--2025 outcomes, and
sharing its A1 arm with the contrasts it prices (its non-circular
content is that the protected twin sides with the retaining arms
exactly when the deficit is positive). The weak/strong scoring
taxonomy below was formalized only after the withheld-era runs, and
the rule as lodged was scoped to walk-forward designs; we decline to
hide behind that scoping and count the stitched designs against
ourselves. One further scoring choice is confessed rather than buried:
no registration fixed whether the gross or the net-of-cost register
adjudicates the NBER cell's strong form---gross, its retaining chain
is partial; net (the economically primary register by the cost
battery's own argument, an argument that predates the cell), it is a
full hit---and the amendment lodged for the constituent test fixes
the net register in advance so the ambiguity cannot recur.

\textbf{The inclusive scorecard.} Scoring every cell the rule could
touch---six withheld-era cells, the two schedule-resampled (stitched)
1990--2025 designs, and the four NBER cells---the record is
\textbf{8/12 under the strong form} (sign predicts the arm-table
direction: ordering when significantly positive, inversion when
negative) and \textbf{12/12 under the weak form}, which is
one-directional: dissociation \emph{only if} $\G$ is significantly
positive---no cell dissociates without a positive deficit; the converse
fails in the two stitched cells, and those failures are counted above
as strong-form misses. The four strong-form misses share one
shape: a marginally significant $\G$ with a flat arm table. Dependence
caveats at full prominence: the withheld 15\% full-era cell is the
union of its two sub-era cells; the four threshold cells share one
history's crises; the stitched designs resample the same panels. The
scorecard's evidential force is closer to two-to-three effectively
independent observations, in both directions, than to twelve
trials---and the two 1990--2025 walk-forward tests that motivated the
rule are not counted at all.

\textbf{Custody.} Because a rule of this shape is exactly the kind of
thing a file drawer manufactures, every registration in the program is
graded by the strength of its timestamp custody, against ourselves
(Table~\ref{tab:custody}). The early batteries are self-attested; the
later batteries have git-separated custody (specification pushed
publicly before results existed); the NBER cell resolved under
third-party (OSF) custody; and all future registrations are OSF-first,
which is the only available answer to the unreportable-frozen-specs
question. As of this record every \emph{run} registration in the
program is closed; open are the two market-clock forward tests---the
temporal-deployment scoring (2027-04-07) and the CRSP constituent
battery---and two lodged-but-unrun registrations we disclose at the
prominence a file drawer would deny them: a granularity-window map and
a regional battery with region-specific sign predictions, scheduled
before any further reanalysis of the 1990--2025 panel.

\begin{table}[t]
\centering
\caption{Provenance grades for the program's registrations (compact;
the full ledger, with commit hashes and verification instructions, is
released). OSF project ID anonymized for review. \todo{replace OSF ID
with anonymized view-only link}}
\label{tab:custody}
\small
\begin{tabular}{@{}p{0.48\columnwidth}p{0.40\columnwidth}@{}}
\toprule
\textbf{Registration} & \textbf{Custody grade} \\
\midrule
Gates + dissociation; drawdown labeler; threshold sweep (2026-07-14) &
self-attested (spec and results entered the repository together) \\
Costs, replay, seed parity; mechanism probes; withheld era (2026-07-15);
banded-baseline control with lodged adverse branch (2026-07-16)
& git-separated (spec pushed before runs; server timestamps) \\
NBER announce-lagged forward test & \emph{resolved under third-party
custody} (OSF 22:47 ET; data joined 22:55 ET; hit) \\
CRSP constituents + T-split amendment; temporal-deployment test &
forward, third-party custody (open) \\
\bottomrule
\end{tabular}
\end{table}

\textbf{The live test.} The rule's standing changes register only with
temporally split measurement, so the program's next cells are already
committed. (1)~\emph{The deployment-form temporal test}: all
diagnostics frozen \textbf{2026-06-30}, on the vendored panel through
2026-05-29
($\G = +0.00081 \pm 0.00033$ on the NBER cell;
$+0.00091 \pm 0.00020$ on the drawdown cell---both positive, so the
lodged prediction is that the retaining arms outperform in the outcome
window), scored against realized reactivations in 2026-07-01 through
2027-03-31 on a pre-named date, \textbf{2027-04-07}, with a void
declared (not a hit) if no qualifying reactivation occurs. One clean
refuting cell falsifies the deployment form of the rule.
(2)~\emph{Constituent-level replication}: the CRSP S\&P~500 panel,
registration lodged before any data access, with a T-split amendment
fixing the net-cost register and the window estimator in advance---the
test that decides whether the deficit survives disaggregation, the
canonical worry for industry-level effects
\cite{moskowitz1999industries,grundy2001understanding}. Both are
reported at headline prominence regardless of outcome.

%=============================================================================
\section{Field kit and limitations}
\label{sec:kit}
%=============================================================================

\subsection{Running the protocol on your own book}

\begin{enumerate}
  \item \textbf{Freeze the design.} Fix the labeler family, thresholds,
  probe length, dormancy cutoff, and test family in writing; lodge the
  document externally (OSF or equivalent) before any run. The record
  above is the argument that this step is load-bearing.
  \item \textbf{Label causally.} Regime labels must be computable from
  information through $t-1$---announcement dates, not revision dates,
  for any official indicator.
  \item \textbf{Run the screen.} One model per regime versus the same
  machinery on $\ge 50$ block-shuffled label draws, on \emph{your
  decision metric}, walk-forward, with a split-half check. A fail
  predicts your regime-aware evaluation should go flat; run it anyway
  and distrust everything if it does not.
  \item \textbf{Measure $\G$.} Two arms: your deployed allocation
  policy, and its twin with per-regime parameters pinned through
  dormancy. $\G$ = probe-window regret difference. Charge your own
  crisis-window cost schedule (both-pay); keep $\G_{\mathrm{net}}$.
  \item \textbf{Price it.} Apply the boxed rule of
  Section~\ref{sec:breakeven} with your sleeve size, measured dormancy
  scale, and honest carrying cost.
  \item \textbf{Decide, and re-estimate.} $\G_{\mathrm{net}}$
  significantly positive and the rule satisfied: retention pays, and
  any isolation-preserving implementation collects it. Otherwise: do
  not buy retention machinery this era---and re-run the diagnostic when
  crisis density shifts, because the paying window's location moved
  across eras in our record.
\end{enumerate}

\subsection{Limitations, register, and what we do not claim}
\label{sec:limits}

\textbf{Single histories.} Every real-data cell is one realized market
history; seed intervals quantify implementation noise conditional on
that history. The record spans 1926--2025 and scores in both
directions, but it is three walk-forward windows and one country.

\textbf{Aggregation.} The equity substrate is industry portfolios;
aggregation itself can manufacture slow dynamics \cite{lo1988stock},
crisis-conditional behavior is a documented feature of this territory
\cite{cooper2004market,daniel2016momentum,henkel2011time,
rapach2010out}, and industry-level effects need not survive
disaggregation \cite{grundy2001understanding}. The lodged CRSP
constituent test is the decider, and until it resolves, the deficit is
a fact about the aggregate cross-section.

\textbf{One labeler family, one era-dependent window.} The positive
deficits live under drawdown- and recession-type labelers with slow
crisis dynamics; the paying threshold relocated across eras. The
protocol ships with the instruction to estimate the window locally,
not with a universal constant.

\textbf{Mechanism.} On real data the measured deficit is a persistent
allocation-level service gap, not demonstrated forgetting; the
controlled dissociation of mechanisms lives in the companion
\cite{companion2026}, and buying our diagnostic does not require
buying any mechanism story.

\textbf{No alpha.} Every number here is decision regret at portfolio
level, gross except where a modeled crisis-window cost schedule is
stated; nothing is a claim of live trading profitability
\cite{mclean2016does}. The protocol's product is a decision---pay to
remember, or don't---and its warranty is the record of Table~
\ref{tab:master}: three flats, two inversions, two orderings, all
published under the same rule.

\bibliographystyle{ACM-Reference-Format}
\bibliography{references_icaif}

\end{document}
